\section*{Chapter \ref{ch:qm:robust-statistics-and-heavy-tails} --- Robust Statistics and Heavy Tails}

\begin{solution}{exo:qm:robust-statistics-and-heavy-tails:1}
10\% for the \omterm{def:qm:robust-statistics-and-heavy-tails:winsor}{trimmed mean} (more than the trimmed share on one side carries it away); 25\% for the interquartile range, which depends only
on the quartiles.
\end{solution}

\begin{solution}{exo:qm:robust-statistics-and-heavy-tails:2}
Moments of order below 3 exist: the mean and the variance. The kurtosis, a fourth moment, does not; the third moment is the borderline case
and is infinite for an exact Pareto tail, since $\int^\infty x^2 \cdot x^{-4}\,x\,dx$ diverges logarithmically.
\end{solution}

\begin{solution}{exo:qm:robust-statistics-and-heavy-tails:3}
$\rho_S = \frac6\pi\arcsin0.21 = 0.404$; $\tau = \frac2\pi\arcsin0.42 = 0.276$.
\end{solution}

\begin{solution}{exo:qm:robust-statistics-and-heavy-tails:4}
$\E[\psi_c'(Z)] = 2\Phi(c) - 1 = 0.8214$. $\E[\psi_c(Z)^2] = \E[Z^2; |Z| \le c] + c^2\P(|Z| > c) = (2\Phi(c) - 1) - 2c\varphi(c) + 2c^2(1 - \Phi(c)) =
0.7102$. The asymptotic variance is $\E[\psi^2]/\E[\psi']^2$, so the efficiency relative to the mean is $0.8214^2/0.7102 = 0.950$.
\end{solution}

\begin{solution}{exo:qm:robust-statistics-and-heavy-tails:5}
Above $x_0 = X_{(k+1)}$ the Pareto density is $\alpha x_0^\alpha x^{-\alpha-1}$. The log-likelihood of the $k$ largest is $k\ln\alpha + k\alpha\ln x_0 - (\alpha +
1)\sum_{i \le k}\ln X_{(i)}$; its derivative $k/\alpha - \sum_i(\ln X_{(i)} - \ln x_0)$ vanishes at $\hat\alpha = k/\sum_i\ln(X_{(i)}/x_0)$, the \omterm{def:qm:robust-statistics-and-heavy-tails:hill}{Hill
estimator}.
\end{solution}

\begin{solution}{exo:qm:robust-statistics-and-heavy-tails:6}
$0.93 + \frac{0.34}{0.069}\bigl((0.05/0.0001)^{0.069} - 1\bigr) = 3.59\%$ (with the unrounded fit; $N_u/n = 355/7\,098 = 0.05$). Two days in 27 years
lost more (3.7\% and 4.7\%), about what the fit implies (0.7 expected); so few observations are why the extrapolation's
uncertainty is large.
\end{solution}

\begin{solution}{exo:qm:robust-statistics-and-heavy-tails:7}
Above the 90\%, 95\%, 97.5\% and 99\% quantiles (710, 355, 178 and 71 excesses): $\xi = 0.033$, $0.069$, $0.098$, $0.212$; one-in-a-thousand-days
loss 2.45\%, 2.47\%, 2.49\%, 2.43\%. The shape wanders, the quantile inside the data does not.
\end{solution}

\begin{solution}{exo:qm:robust-statistics-and-heavy-tails:8}
With 250 days, winsorising at 1\% caps the two or three largest moves on each side. It happens to neutralise the chapter's single bad print
(the corrupted standard deviation becomes 0.341\% against a clean 0.339\%), but it also caps genuine large moves every day, biasing the
clean volatility down (0.328\% against 0.339\%), most in the stress periods that matter; and three bad prints in a year pass through. The
standard deviation of winsorised data also needs a consistency correction. Filter prints upstream, and use a robust scale by design.
\end{solution}

\begin{solution}{pb:qm:robust-statistics-and-heavy-tails:1}
\textbf{1.} $\ln(1.5172/1.1572) = 0.271$: the returns of 24 and 25 March 2026 become about $+27\%$ and $-27\%$.
\textbf{2.} Clean: 0.339\%, 0.275\%, 0.300\%. Corrupted: 2.434\%, 0.279\%, 0.305\%.
\textbf{3.} The mean ($-0.011\%$) and the median ($-0.030\%$) are unchanged, the first by luck, since the two bad returns cancel; the Huber
location barely moves ($-0.01953\%$ to $-0.01948\%$), its bounded influence at work.
\textbf{4.} Seventeen moves beyond four standard deviations (six losses, eleven gains), against 0.45 expected under a normal law.
\textbf{5.} 6.3; a fourth moment of a law whose \omterm{def:qm:robust-statistics-and-heavy-tails:heavy}{tail index} is near four is barely finite, and a handful of days decide the estimate.
\textbf{6.} About 3.9 on a plateau for $k$ between 100 and 200 largest losses (4.9 at $k = 50$, drifting to 2.6 at $k = 700$).
\textbf{7.} $u = 0.93\%$, 355 excesses, $\xi = 0.069$, $\beta = 0.34\%$.
\textbf{8.} $\nu = 4.7$, scale 0.44\%.
\textbf{9.} Normal 1.80\%, Student $t$ 2.74\%, GPD 2.47\%; empirical 2.27\%.
\textbf{10.} 35 days beyond the normal quantile, 6 beyond the GPD quantile.
\textbf{11.} Normal 2.27\% (+27\%), Student $t$ 2.91\% (+6\%), GPD 2.85\% (+15\%).
\textbf{12.} The normal quantile is the standard deviation times 3.09, and the standard deviation has an unbounded \omterm{def:qm:robust-statistics-and-heavy-tails:influence}{influence function}; the pair
of $\pm 27\%$ returns raises it from 0.58\% to 0.74\%.
\textbf{13.} It rises from 0.07 to 0.22: the bad loss is by far the largest excess, and the shape parameter is fitted mostly by the largest
excesses.
\textbf{14.} Its likelihood weights an observation by $(\nu + 1)/(\nu + z^2)$, which falls with the size of the residual: the $t$ likelihood is itself a
robust \omterm{def:qm:estimation:m}{M-estimator}.
\textbf{15.} Partly: a MAD-based normal quantile ignores the bad print but is even further from the true tail (the normal shape is wrong); robust
scale does not make a thin-tailed model fit a \omterm{def:qm:robust-statistics-and-heavy-tails:heavy}{heavy-tailed} variable.
\textbf{16.} The GPD or Student-$t$ quantile, about 2.5\% to 2.7\%, reported with its threshold and sensitivity, never the normal 1.8\%.
\textbf{17.} At ingestion: a filter that flags a print whose return exceeds many robust scales and reverses the next day, before any statistic
is computed.
\textbf{18.} Joint losses are more frequent than a Gaussian \omterm{def:qm:robust-statistics-and-heavy-tails:copula}{copula} implies (32\% against 20\% in the joint 5\% tail), so a limit on the sum of the
two exposures computed with correlation 0.42 understates their joint stress loss.
\textbf{19.} Named result: \emph{the bad tick}: the one-in-a-thousand-days EUR/USD loss is 1.80\% (normal), 2.74\% (Student $t$) and 2.47\%
(GPD above the 95\% quantile); one transposed-digit print raises them by 27\%, 6\% and 15\%.
\textbf{20.} Estimates of the centre and the scale of the bulk can be made robust; estimates of the tail cannot, because in the tail every
observation is information, so the protection must come from cleaning the data.
\end{solution}

\begin{solution}{iq:qm:robust-statistics-and-heavy-tails:1}
The data: a bad print or a stale-then-corrected price produces a pair of large opposite returns. Look at the largest returns in the window,
compare with another source, and compare the standard deviation with the MAD: a sevenfold gap between them is a data problem, not a market.
\iqlookfor{Data before model; the paired-reversal signature; a robust scale as a diagnostic.}
\end{solution}

\begin{solution}{iq:qm:robust-statistics-and-heavy-tails:2}
The smallest fraction of arbitrary contamination that can carry the estimate arbitrarily far. Mean: $1/n$, tending to zero. Median: one half.
MAD: one half.
\iqlookfor{The definition, and the contrast between zero and one half.}
\end{solution}

\begin{solution}{iq:qm:robust-statistics-and-heavy-tails:3}
Moments of order below 3 exist and above 3 do not: finite variance, infinite kurtosis. The sample kurtosis then does not converge; it grows with
the sample and is dominated by the largest observation, so it should not be used as a parameter.
\iqlookfor{The moment rule, and the consequence for sample moments.}
\end{solution}

\begin{solution}{iq:qm:robust-statistics-and-heavy-tails:4}
About six observations lie beyond the 99.9\% quantile in 25 years: the empirical quantile is noisy, a normal fit is badly biased. Fit a
generalised Pareto law above a high threshold (\omterm{def:qm:robust-statistics-and-heavy-tails:gpd}{peaks over threshold}), check stability across thresholds, compare with a Student-$t$ fit,
handle volatility clustering by fitting to returns scaled by a volatility forecast, and report the uncertainty.
\iqlookfor{\omterm{def:qm:robust-statistics-and-heavy-tails:gev}{Extreme-value theory} with threshold diagnostics, and awareness of dependence.}
\end{solution}

\begin{solution}{iq:qm:robust-statistics-and-heavy-tails:5}
\omterm{def:qm:robust-statistics-and-heavy-tails:rank}{Rank correlations} are invariant to increasing transformations and bounded in influence, so outliers and heavy tails do not dominate them. Neither
Pearson nor \omterm{def:qm:robust-statistics-and-heavy-tails:rank}{rank correlation} describes the tails: two pairs with the same \omterm{def:qm:robust-statistics-and-heavy-tails:rank}{rank correlation} can have very different joint-crash
probabilities.
\iqlookfor{Invariance and robustness of ranks, and tail dependence as a separate quantity.}
\end{solution}

\begin{solution}{iq:qm:robust-statistics-and-heavy-tails:6}
Any joint law is its margins composed with a \omterm{def:qm:robust-statistics-and-heavy-tails:copula}{copula}, unique for continuous margins, so dependence can be modelled separately from the margins. The
Gaussian \omterm{def:qm:robust-statistics-and-heavy-tails:copula}{copula} has zero tail dependence for any correlation below one: extreme joint moves become asymptotically independent, while in markets
they cluster, as a Student-$t$ \omterm{def:qm:robust-statistics-and-heavy-tails:copula}{copula} allows.
\iqlookfor{The decomposition, and tail dependence zero versus positive.}
\end{solution}
